% 样本压缩、信息依赖与算法稳定性的三种比较实验（示意，跨栏约 165 mm）。
% 左：6 条训练记录中保留 2 条（绿圈）及旁信息 u，重构 rho(S_I,u)；其余 4 条检验重构结果。
% 中：二元观测 z 与输出 W 的联合分布（对角 0.45、非对角 0.05）与乘积分布（每格 0.25），
%     二者的相对熵即互信息 log2-h(0.1)≈0.368，与正文图 fig:lt7a-joint-product-information 相同。
% 右：相邻样本 S 与 S^{(i)} 只在第 3 项不同（红），两次输出在同一观测 z 上的损失差。
\begingroup
\providecolor{FigureGreen}{HTML}{829C8B}
\providecolor{FigureGreenFill}{HTML}{E8EEE9}
\begin{tikzpicture}[figure picture,x=1mm,y=1mm]
\tikzset{dot/.style={circle,fill=FigureBlue,inner sep=0pt,minimum size=2.6mm},
 odot/.style={circle,draw=FigureBlue,fill=white,line width=.7pt,inner sep=0pt,minimum size=2.6mm},
 box/.style={rectangle,draw=FigureStroke,fill=white,line width=.6pt,rounded corners=1mm,inner sep=2pt,font=\figurefont\footnotesize,align=center},
 cell/.style={rectangle,draw=FigureStroke,line width=.5pt,minimum width=8mm,minimum height=7mm,inner sep=0pt,font=\figurefont\footnotesize}}
% ---------- 面板 1：样本压缩
\figurePanel{0}{0}{52}{60}{样本压缩：能否重构}{FigureBlueFill}
\foreach \x/\t in {6/o,13/d,20/o,27/d,34/o,41/o}{\ifx d\t\node[dot] at (\x,44) {};\else\node[odot] at (\x,44) {};\fi}
\foreach \x in {13,27}{\draw[FigureGreen,line width=1pt] (\x,44) circle (2.6);}
\node[figure note,anchor=west] at (44,44) {$S$};
\node[box] (u) at (46,35) {$u$};
\node[box] (rec) at (26,25) {重构 $\varrho(S_I,u)$};
\draw[figure flow] (13,41)--(rec.north west);
\draw[figure flow] (27,41)--(rec.north);
\draw[figure flow] (u.south)--(rec.north east);
\node[figure note] at (26,13) {其余记录检验重构结果\\控制：描述个数\\$2^b\sum_{j\leq k}\binom nj$};
% ---------- 面板 2：信息依赖
\figurePanel{58}{0}{52}{60}{信息依赖：透露多少}{FigureYellowFill}
\node[figure note] at (71,47.5) {联合 $P_{S,W}$};
\node[figure note] at (97,47.5) {乘积 $P_S\otimes P_W$};
\node[cell,fill=FigureYellowFill] at (67,40) {$.45$}; \node[cell,fill=white] at (75,40) {$.05$};
\node[cell,fill=white] at (67,33) {$.05$}; \node[cell,fill=FigureYellowFill] at (75,33) {$.45$};
\foreach \x in {93,101}{\foreach \y in {33,40}{\node[cell,fill=FigurePaper] at (\x,\y) {$.25$};}}
\node[font=\figurefont\footnotesize] at (84,36.5) {$\|$};
\node[figure note] at (84,25) {$D_{\mathrm{KL}}=I(W;S)\approx0.368$};
\node[figure note] at (84,13) {参照：输出与样本独立\\控制：互信息、条件互信息\\或后验—先验相对熵};
% ---------- 面板 3：算法稳定性
\figurePanel{116}{0}{52}{60}{算法稳定性：换一项变多少}{FigureRedFill}
\foreach \x in {126,132,144,150}{\node[dot] at (\x,44) {}; \node[dot] at (\x,34) {};}
\node[dot] at (138,44) {};
\node[circle,fill=FigureRed,inner sep=0pt,minimum size=2.6mm] at (138,34) {};
\draw[FigureRed,line width=.6pt,dash pattern=on 1.5pt off 1pt] (138,41.4)--(138,36.6);
\node[figure note,anchor=east] at (123.5,44) {$S$};
\node[figure note,anchor=east] at (123.5,34) {$S^{(i)}$};
\node[box,anchor=west] (a1) at (154.5,44) {$A(S)$};
\node[box,anchor=west] (a2) at (154.5,34) {$A(S^{(i)})$};
\draw[figure flow] (151.8,44)--(a1.west);
\draw[figure flow] (151.8,34)--(a2.west);
\node[figure note] at (142,25) {$|\ell(A(S),z)-\ell(A(S^{(i)}),z)|\leq\beta_n$};
\node[figure note] at (142,13) {参照：相邻训练集\\控制：损失差、参数距离\\或输出分布的距离};
\end{tikzpicture}
\endgroup
